\subsection{בחירת ה־\textenglish{Baseline}}

בפרויקט כמו \textenglish{AutoScout}, שבו ההחלטה הסופית מתקבלת מצינור
רב־שלבי, ה־\textenglish{Baseline} המתאים הוא גרסה פשוטה יותר של אותו
צינור ולאו דווקא מודל חלופי. מטרת ההשוואה היא לבדוק מה כל שכבה מוסיפה:
צמצום מועמדים לפי המקצה, צבע ברית, זיהוי ספרות, מעקב וצבירת ראיות
לאורך זמן.

\subsection{קווי הבסיס שנבחרו}

להערכת המערכת ובידוד התרומה של כל שכבת לוגיקה, הוגדרו שני סוגים שונים
של קווי בסיס:

\paragraph{קווי בסיס תיאורטיים-הסתברותייים:}
קווי בסיס אלו מייצגים מערכת ``עיוורת'' לחלוטין ומגדירים את רצפת
הביצועים הסטטיסטית התיאורטית של הבעיה:

\begin{itemize}
\item \textbf{ניחוש אקראי מתוך שש קבוצות:} מאחר שבכל מקצה משתתפות שש
  קבוצות בלבד, מערכת שאין לה מידע נוסף יכולה רק לנחש. הדיוק הצפוי הוא
  \textbf{\(1/6\)}, כלומר כ־\textbf{0.17}.
  
\item \textbf{ניחוש מתוך הברית הנכונה בלבד:} אם ידוע רק צבע הברית,
  מרחב החיפוש מצטמצם לשלוש קבוצות. במקרה כזה הדיוק הצפוי הוא
  \textbf{\(1/3\)}, כלומר כ־\textbf{0.33}.
\end{itemize}

\paragraph{קו בסיס אלגוריתמי־אמפירי:}

\begin{itemize}
\item \textbf{זיהוי ספרות נקודתי ללא צבירה לאורך זמן:} קו בסיס זה
  משתמש בזיהוי הספרות שעל הבאמפר ומשווה אותן לקבוצות האפשריות, אך מקבל
  החלטה על סמך פריים בודד. הוא טוב יותר מניחוש, אבל רגיש לטשטוש תנועה,
  לחסימות ולזיהוי חלקי. קו בסיס זה נמדד אמפירית על נתוני האמת, ומטרתו
  להוכיח את נחיצותה של שכבת קבלת ההחלטות במערכתית שפיתחתי.
\end{itemize}

\subsection{ניתוח ההשוואה}

מאחר שקווי הבסיס התיאורטיים מייצגים גבולות הסתברותיים מתמטיים ולא
מערכות תוכנה שרצו בפועל, הטבלה שלהלן מפרידה בין המדדים התיאורטיים לבין
תוצאות הניסוי האמפירי.

הניסוי האמפירי בוצע באמצעות תהליך תיקוף ידני מבוסס דגימה, על מקצי
תחרות מלאים, כאשר בשל אילוצי תיוג ידני, קו הבסיס האלגוריתמי נמדד על
תת־קבוצה מייצגת של שני מקצים:

\begin{table}[H]
  \centering
  \begin{tabular}{|>{\raggedleft\arraybackslash}p{0.125\textwidth}|>{\raggedleft\arraybackslash}p{0.125\textwidth}|>{\raggedleft\arraybackslash}p{0.125\textwidth}|>{\raggedleft\arraybackslash}p{0.125\textwidth}|>{\raggedleft\arraybackslash}p{0.125\textwidth}|>{\raggedleft\arraybackslash}p{0.125\textwidth}|>{\raggedleft\arraybackslash}p{0.125\textwidth}|}
    \hline
    \textbf{שיטה} & \textbf{סוג המדד} & \textbf{מספר מקצים שנמדדו} & \textbf{סך מופעי רובוט} & \textbf{סך שיוכים נכונים} & \textbf{תוצאה} \\ \hline
    ניחוש אקראי & תיאורטי & -- & -- & -- & \textbf{0.17} \\ \hline
    ניחוש לפי ברית & תיאורטי & -- & -- & -- & \textbf{0.33} \\ \hline
    זיהוי ספרות ללא צבירה & אמפירי & 2 & 183 & 102 & \textbf{0.56} \\ \hline
    המערכת המלאה & אמפירי & 3 & 273 & 204 & \textbf{0.75} \\ \hline
  \end{tabular}
  \caption{השוואת קווי בסיס למערכת המלאה.}
  \label{tab:baseline_comparison}
\end{table}

ההשוואה מראה שכל רכיב מוסיף סוג אחר של מידע:
 
\begin{itemize}
\item \textbf{מספר המקצה} מצמצם מיד את מרחב החיפוש לשש קבוצות בלבד.
\item \textbf{צבע הברית} מצמצם אותו עוד יותר לשלוש מועמדות אפשריות.
\item \textbf{זיהוי הספרות} מוסיף רמז זהותי חזק, אך הוא רועש כאשר הבאמפר
  מוסתר, מטושטש או מצולם בזווית קשה.
\item \textbf{המעקב וההצבעה המצטברת} מאפשרים לצבור ראיות לאורך זמן,
  ולכן מפחיתים את ההשפעה של טעויות מקומיות בפריימים בודדים.
\end{itemize}

מכאן שהיתרון של המערכת המלאה אינו נובע רק ממודל ספרות טוב יותר, אלא
מכך שהיא פותרת את הבעיה ברמת \textbf{עקבה שלמה} ולא ברמת פריים יחיד.

\subsection{השוואה לפתרונות קיימים}

אפשר להשוות את הפרויקט גם לשתי משפחות נפוצות של פתרונות:

\begin{itemize}
\item \textbf{סקאוטינג ידני:} גמיש ועתיר הקשר אנושי, אך דורש כוח אדם רב,
  אינו שומר מסלולים מלאים, ורגיש לעייפות ולטעויות הקלדה.

\item \textbf{מערכות \textenglish{Detection + Tracking} כלליות:} מסוגלות
  לזהות ולעקוב אחר רובוטים, אך ללא ידע דומייני נוסף הן אינן פותרות היטב
  את שאלת הזהות של כל רובוט כאשר כמה רובוטים דומים מופיעים יחד על
  המגרש.
\end{itemize}

היתרון של \textenglish{AutoScout} הוא בשילוב בין אוטומציה לבין ידע
דומייני: מצד אחד המערכת מהירה ועקבית יותר מסקאוטינג ידני, ומצד שני היא
מנצלת את לוח המשחקים של המקצה --- מידע שמערכות מעקב כלליות אינן מנצלות.

\subsection{מסקנה מן ההשוואה}

המסקנה המרכזית היא שהפתרון החזק ביותר לבעיה זו אינו מודל יחיד, אלא
מערכת שמחברת זיהוי, מעקב וידע מוקדם. זו גם הנקודה שבה \textenglish{AutoScout}
שונה מקווי בסיס פשוטים יותר ומכלי מעקב כלליים.
